\section{Discussion and conclusions}
\label{sec:conclusions}

The principal result is the sixteen-cell cross section of Sec.~\ref{sec:result16}, with its covariance and generator
comparison.

\label{sec:limits}
Two validation findings point to the unfolded result retaining part of its prior.
The first is the residual under a generator-sized change of the truth in Sec.~\ref{sec:departure}, which we
interpret as regularization bias.
The second is the envelope of Sec.~\ref{sec:uncertainty-status}, which tracks the same bias and alone exceeds the
precision target in two of the three tested reporting definitions.
The note's calibrated tests show the same pull: their null distributions for the external predictions are displaced
upward by it, as the calibration reproduces (\nsec{sec:joint5d-results}).
The miscoverage of Sec.~\ref{sec:coverage} is also attributed to a residual bias, but whether that bias has the same
origin is not established.

The studies of Sec.~\ref{sec:validation} show what must be checked before unfolded intervals are interpreted:
closure with the backgrounds included; truth changes of generator size, not only nominal closure; coverage under
changes of the truth, not only at nominal truth; stability under training seeds and rounding-level input changes;
and run-period consistency.

\subsection{Limitations}
\label{sec:limitations}

The limitations of these results are stated with the results they qualify, in
Secs.~\ref{sec:method}--\ref{sec:result16}, and those of the calibrated tests in the companion note.
The low stability of four of the sixteen cells is stated only here.
The principal ones are these.
\begin{enumerate}
\item The two-dimensional covariance has not been independently validated, and its coverage has not been tested
(Sec.~\ref{sec:uq2d}).
\item The three- and four-observable central values lack intervals of validated coverage (Sec.~\ref{sec:central}).
\item The sixteen-cell result also has total uncertainties, whose coverage is not validated; Sec.~\ref{sec:fd16}
reports the fake-data tests of the estimator's bias.
\item Four of the sixteen cells, \fvval{h2-lowstab-cells}, have a stability of
\fvval{h2-lowstab-min}--\fvval{h2-lowstab-max}: of the simulated signal events generated in such a cell and
reconstructed, only that fraction is reconstructed in the same cell.
\item The note's calibrated tests reject hybrid nulls, conditionally on the simulation's hadronic response.
However, the rejections carry over to the generator predictions themselves only under a bound on the residual below
the fine grid, which has not been shown.
Their calibration is valid to the extent that the simulated chain reproduces the analysis of the data
(\nsec{sec:joint5d-conditions}).
They are not physical exclusions of the generators, and they show no evidence of additional discrimination beyond
matched coarse projections.
\end{enumerate}
